% ============================================================
% Capitolo 9 — Confidence Tags
% ============================================================

\chapter{Confidence Tag: FACT, INFERRED, ASSUMPTION}
\label{ch:confidence-tags}

Gli agenti AI hanno un problema strutturale: rispondono con la stessa sicurezza apparente sia quando sanno, sia quando deducono, sia quando inventano. I confidence tag sono il meccanismo con cui ADLC rende esplicita questa distinzione.

\section{I tre tag}
\label{sec:09-tre-tag}

I tag seguono sempre lo stesso formato, in coda all'affermazione che classificano:

\begin{lstlisting}[language=,caption={Formato dei confidence tag}]
AI CONFIDENCE: [FACT | INFERRED | ASSUMPTION]
Basis: [motivazione]
[TODO: azione suggerita — obbligatorio per ASSUMPTION su HIGH-risk]
\end{lstlisting}

\subsection*{FACT}

L'agente ha verificato l'informazione leggendo direttamente il codice, un file, o un output di test presente nel contesto.

\begin{lstlisting}[language=,caption={Esempio FACT}]
La funzione getUserById lancia UserNotFoundError su record mancante.

AI CONFIDENCE: FACT
Basis: letto src/services/userService.ts, riga 45-50
\end{lstlisting}

\subsection*{INFERRED}

L'agente ha dedotto l'informazione da contesto indiretto: nome della funzione, struttura del codice, convenzioni del framework. La deduzione è logica ma non verificata direttamente.

\begin{lstlisting}[language=,caption={Esempio INFERRED}]
Il middleware usa @fastify/jwt basandosi sull'import in src/app.ts.

AI CONFIDENCE: INFERRED
Basis: import @fastify/jwt in src/app.ts; non ho letto la
       configurazione completa del middleware
\end{lstlisting}

\subsection*{ASSUMPTION}

L'agente ha fatto un'ipotesi non supportata da evidenza diretta nel contesto.

\begin{lstlisting}[language=,caption={Esempio ASSUMPTION}]
Ho assunto che tasks abbia un indice su user_id (PERF-01: P95 < 200ms).

AI CONFIDENCE: ASSUMPTION
Basis: non ho letto lo schema Prisma
TODO: verifica index su prisma/schema.prisma prima di andare in prod
\end{lstlisting}

\section{Quando i tag sono obbligatori}
\label{sec:09-quando-obbligatori}

\textbf{Sempre obbligatori:}
\begin{itemize}
  \item Output classificati HIGH-risk o superiore
  \item Claim su vincoli SEC-XX o PERF-XX
  \item Codice che modifica path presenti negli HALT trigger
  \item Affermazioni su comportamento del codice esistente senza averlo letto
\end{itemize}

\textbf{Consigliati ma non obbligatori in Mode LITE:}
\begin{itemize}
  \item Output MEDIUM superiori a 10 righe
  \item Deduzioni su comportamento di librerie esterne
\end{itemize}

\section{Il tag come segnale di azione}
\label{sec:09-segnale}

\begin{table}[htbp]
  \centering
  \begin{tabular}{ll}
    \toprule
    \textbf{Tag} & \textbf{Azione} \\
    \midrule
    FACT & Usa l'informazione. Verifica solo se la fonte potrebbe essere obsoleta. \\
    INFERRED & Verifica prima di agire su operazioni MEDIUM o superiori. \\
    ASSUMPTION & \textbf{Stop.} Verifica prima di qualsiasi operazione. \\
    \bottomrule
  \end{tabular}
  \caption{Azione da intraprendere per ogni confidence tag}
  \label{tab:09-azioni}
\end{table}

\section{In pratica su TaskFlow API}
\label{sec:09-pratica}

Lorenzo chiede all'agente di rivedere la gestione errori di \texttt{POST /tasks}. L'agente risponde con tre claim di certezza diversa:

\begin{itemize}
  \item \textbf{FACT}: la validazione Zod è corretta --- verificato leggendo il file.
  \item \textbf{INFERRED}: il middleware è probabilmente registrato nell'ordine giusto --- consigliato di verificare.
  \item \textbf{ASSUMPTION}: il mapping degli errori Prisma potrebbe non esistere --- Lorenzo deve verificare prima di andare in staging.
\end{itemize}

Senza i confidence tag, Lorenzo avrebbe ricevuto lo stesso testo senza sapere cosa è stato verificato e cosa no.

\section*{\LabelSummary}

\begin{itemize}
  \item Tre tag: FACT (verificato), INFERRED (dedotto), ASSUMPTION (ipotesi).
  \item FACT: usa. INFERRED: verifica prima di operazioni MEDIUM+. ASSUMPTION: stop, verifica subito.
  \item Obbligatori su output HIGH-risk, claim SEC/PERF, codice in zone HALT.
  \item In Mode LITE si riducono ai soli casi critici.
\end{itemize}
